\documentclass[11pt]{article}
\usepackage[utf8]{inputenc}
\usepackage[spanish,es-noindentfirst]{babel}
\usepackage[a4paper,margin=2cm]{geometry}
\usepackage[table]{xcolor}
\usepackage{array}
\usepackage{longtable}
\usepackage{enumitem}
\usepackage{fancyhdr}
\usepackage{hyperref}

\definecolor{vino}{HTML}{621F32}
\definecolor{oro}{HTML}{BC955C}
\definecolor{grisclaro}{HTML}{F1F5F9}
\definecolor{grisborde}{HTML}{CBD5E1}

\hypersetup{
  pdftitle={Caso de Prueba F27 - Torre Caballito 3D},
  pdfauthor={Eje Central - QA},
  colorlinks=true,
  linkcolor=vino,
  pdfborder={0 0 0}
}

\pagestyle{fancy}
\fancyhf{}
\renewcommand{\headrulewidth}{0.6pt}
\renewcommand{\footrulewidth}{0.3pt}
\fancyhead[L]{\small\color{vino}\textbf{CP-F27} — Torre Caballito 3D}
\fancyhead[R]{\small\color{grisborde} eje\_central\_front}
\fancyfoot[C]{\small\thepage}

\newcommand{\seccion}[1]{%
  \vspace{6pt}
  {\large\color{vino}\textbf{#1}}\\[2pt]
  {\color{oro}\rule{\linewidth}{1.2pt}}\\[6pt]
}

% Checkbox de PDF interactivo (AcroForm)
\newcommand{\chk}[1]{\CheckBox[name=#1,width=11pt,height=11pt,bordercolor=grisborde,backgroundcolor=1 1 1]{}}

% Campo de texto de una línea
\newcommand{\txtline}[3]{\TextField[name=#1,width=#2,height=13pt,bordercolor=grisborde,backgroundcolor=0.97 0.98 1]{#3}}

\begin{document}
\begin{Form}

% ---------- Portada / encabezado ----------
\begin{center}
  {\Huge\color{vino}\textbf{Caso de Prueba}}\\[4pt]
  {\Large\color{oro}\textbf{CP-F27}}\\[2pt]
  {\large Torre Caballito 3D — Búsqueda de empleado y beacon de ubicación}\\[2pt]
  {\small Módulo: Plantilla de Empleados $\to$ Distribución Geográfica $\to$ Torre Caballito 3D}
\end{center}

\vspace{10pt}

\noindent
\begin{tabular}{|>{\columncolor{grisclaro}}p{3.2cm}|p{12.8cm}|}
\hline
\textbf{ID} & CP-F27-01 \\ \hline
\textbf{Funcionalidad} & F27 — Distribución Geográfica: Torre Caballito 3D (búsqueda de empleado) \\ \hline
\textbf{Página / Ruta} & \texttt{/dashboard/plantilla\_empleados} $\to$ tab \emph{Distribución Geográfica} $\to$ sub-vista \emph{Torre Caballito 3D} \\ \hline
\textbf{Archivo fuente} & \texttt{src/app/dashboard/plantilla\_empleados/\_components/tabs/torre-3d/TorreCaballito3DTab.jsx} \\ \hline
\textbf{Prioridad} & Media \\ \hline
\textbf{Tipo} & Funcional — UI interactiva (3D) \\ \hline
\textbf{Ejecutor} & \txtline{tester_nombre}{9cm}{} \\ \hline
\textbf{Fecha de ejecución} & \txtline{fecha_ejecucion}{5cm}{} \\ \hline
\textbf{Entorno} & \chk{env_local} Local \quad \chk{env_staging} Staging \quad \chk{env_prod} Producción \\ \hline
\end{tabular}

\seccion{Objetivo}
Verificar que el buscador de empleados dentro de la vista 3D de la Torre Caballito
localiza correctamente a un empleado, anima la cámara hasta su piso asignado y
despliega el \emph{beacon} (marcador pulsante + etiqueta flotante) en el piso correcto,
manejando también el caso de un empleado sin piso asignado.

\seccion{Precondiciones}
\begin{itemize}[leftmargin=1.6em]
  \item[\chk{pre1}] Sesión iniciada con permiso de acceso al módulo \emph{Plantilla de Empleados}.
  \item[\chk{pre2}] Se cuenta con el nombre de al menos un empleado \textbf{con} piso asignado en la base (\texttt{piso\_num} no nulo).
  \item[\chk{pre3}] Se cuenta con el nombre de al menos un empleado \textbf{sin} piso asignado (\texttt{piso\_num} nulo), para el caso negativo.
  \item[\chk{pre4}] Navegador de escritorio con soporte WebGL habilitado.
\end{itemize}

\noindent\textbf{Datos de prueba (llenar antes de ejecutar):}\\[4pt]
\begin{tabular}{|p{6.5cm}|p{9.5cm}|}
\hline
Empleado con piso asignado & \txtline{dato_emp_con_piso}{8.8cm}{} \\ \hline
Empleado sin piso asignado & \txtline{dato_emp_sin_piso}{8.8cm}{} \\ \hline
\end{tabular}

\seccion{Pasos de ejecución y resultados esperados}

\begin{longtable}{|c|p{6.7cm}|p{6.7cm}|c|}
\hline
\rowcolor{grisclaro}
\textbf{\#} & \textbf{Acción} & \textbf{Resultado esperado} & \textbf{OK} \\ \hline
\endhead

1 &
Entrar a \emph{Plantilla de Empleados} $\to$ tab \emph{Distribución Geográfica} $\to$ \emph{Torre Caballito 3D}. &
Se muestra spinner "Construyendo Torre Caballito en 3D..." y luego el edificio 3D (32 pisos) con controles de cámara (arrastrar = orbitar). &
\chk{r1} \\ \hline

2 &
Escribir 1 o 2 caracteres del nombre del empleado con piso asignado en el buscador (arriba a la derecha). &
No aparece dropdown de resultados aún (mínimo 3 caracteres); puede verse el spinner de búsqueda. &
\chk{r2} \\ \hline

3 &
Completar al menos 3 caracteres del nombre. &
Tras $\sim$400\,ms aparece dropdown con tarjetas de resultados: nombre, Unidad Administrativa y chip \texttt{Piso N} (o "Piso no asignado"). &
\chk{r3} \\ \hline

4 &
Hacer clic en el resultado del empleado \textbf{con} piso asignado. &
El dropdown se cierra; el input muestra el nombre elegido; la cámara vuela suavemente hacia el piso correspondiente. &
\chk{r4} \\ \hline

5 &
Observar el piso destino tras el vuelo de cámara. &
El piso queda resaltado (borde blanco brillante, tono destacado) y muestra un \emph{beacon}: disco/anillo pulsante azul y una etiqueta flotante con el ícono de pin de ubicación, el nombre del piso y el nombre del empleado. &
\chk{r5} \\ \hline

6 &
Revisar el panel inferior derecho. &
Muestra el nombre del piso, total de activos en el piso y lista de Unidades Administrativas presentes, con botón "Ver todos los empleados del piso". &
\chk{r6} \\ \hline

7 &
Clic en "Ver todos los empleados del piso". &
Se abre modal (\texttt{EmpleadosTableModal}) con el listado de empleados de ese piso; loading visible mientras carga. &
\chk{r7} \\ \hline

8 &
Cerrar el modal y pulsar el botón de "Restablecer Vista" (icono de flecha circular). &
Selección, búsqueda y panel se limpian; la cámara regresa a la posición inicial general. &
\chk{r8} \\ \hline

9 &
Repetir búsqueda usando el nombre del empleado \textbf{sin} piso asignado y seleccionarlo del dropdown. &
Se dispara alerta: "El empleado \{Nombre\} no tiene un piso asignado específico en la base de datos." y \textbf{no} se mueve la cámara ni se selecciona ningún piso. &
\chk{r9} \\ \hline

10 &
Alternar el toggle "Mapa de Calor" $\leftrightarrow$ "Color por Unidad". &
El color de los pisos cambia de escala de calor (amarillo$\to$rojo según densidad) a color por Unidad Administrativa dominante; leyenda lateral cambia acorde. &
\chk{r10} \\ \hline

\end{longtable}

\seccion{Criterios de aceptación}
\begin{itemize}[leftmargin=1.6em]
  \item[\chk{ac1}] La búsqueda no dispara petición al backend con menos de 3 caracteres.
  \item[\chk{ac2}] El resultado seleccionado siempre corresponde al piso real del empleado (\texttt{piso\_num}).
  \item[\chk{ac3}] El caso de empleado sin piso asignado no rompe la vista ni deja el buscador en estado inconsistente (input, dropdown, cámara).
  \item[\chk{ac4}] El botón "Restablecer Vista" limpia sin excepción \emph{todos} los estados: búsqueda, selección de piso, selección de UA, hover.
  \item[\chk{ac5}] No hay errores en la consola del navegador durante todo el flujo.
\end{itemize}

\seccion{Resultado de la ejecución}
\noindent
\chk{veredicto_aprobado} \textbf{Aprobado} \qquad
\chk{veredicto_rechazado} \textbf{Rechazado} \qquad
\chk{veredicto_bloqueado} \textbf{Bloqueado}

\vspace{8pt}
\noindent\textbf{Observaciones / defectos encontrados:}\\[4pt]
\TextField[name=observaciones, width=16cm, height=3cm, multiline=true, bordercolor=grisborde, backgroundcolor=0.97 0.98 1]{}

\vspace{10pt}
\noindent\textbf{Evidencia adjunta (rutas de captura/video):}\\[4pt]
\txtline{evidencia}{16cm}{}

\end{Form}
\end{document}
